% !TeX root = ../教材题库_全部条目_转写工作稿.tex
% Source: XJ-B1-C06-PR-P217-Y1462
% 第6章 / 6.2 抽样 / 教材 p.217 / PDF 第224页
% OCR未校对：公式与文字均可能有误；未完成数学TeX转写；禁止直接发给学生
% 原题及相邻原解截图（核对用）：../题目原图/XJ-B1-C06-PR-P217-Y1462-01.webp
% 本文件只在转写完成后作为题干使用；OCR 条目尚不是可靠题干。
\begin{bankitem}{p.217 练习 第2题}{第6章 / 6.2 抽样 / 教材 p.217 / PDF 第224页}{XJ-B1-C06-PR-P217-Y1462}{OCR未校对：公式与文字均可能有误；未完成数学TeX转写；禁止直接发给学生}
\begin{ocrdraft}
[教材p.217 / PDF第224页]
\par
2．立德中学高一年级有600名学生，为了解这些学生的身高状况，从中抽取
\par
一个容量为60的简单随机样本．试借助Excel或GeoGebra软件生成随机数，完
\par
成这一抽样．
\end{ocrdraft}
\end{bankitem}
